\color{unidarkgreen}{\subsection{Контроль силового выключателя}}
\color{black}

\subsubsection{Общие сведения}

\begin{enumerate}[label=\arabic{section}.\arabic{subsection}.\arabic{subsubsection}.\arabic*,  align=left, labelsep=4pt, leftmargin=0pt, itemindent=57pt]

\item
Функциональный блок контроля силового выключателя (КСВ) предназначен для обработки
технологической информации от выключателя и его цепей управления, а также от трансформатора тока,
позволяя определять их состояние, информировать о возможных неисправностях и формировать
соответствующую сигнализацию.

\item
В случае присоединения линии к системе через два выключателя для каждого из них можно
задействовать свою логику КСВ. Далее приведено описание КСВ для выключателя В1(ЛВ)ВН.

\begin{figure}[H]
\centering
\includegraphics[width=0.8\textwidth,height=0.8\textheight,keepaspectratio]{\fbpath/761.3371 КСВ/img/CTRL_Q1.pdf}
\captionof{figure}{Функциональная схема контроля исправности выключателя}\label{pic:CTRL_Q1}
\end{figure}

\item
Контроль целостности цепей управления осуществляется по сигналам состояния цепей первой и
второй группы электромагнитов отключения (ЭМО1 и ЭМО2) и цепи электромагнита включения (ЭМВ).
Наличие второй группы электромагнитов отключения определяется положением программного
переключателя \textbf{«Кол-во ЭМО В1(ЛВ)ВН»}. При одновременном наличии или отсутствии сигналов контроля цепей ЭМО
и ЭМВ в течение регулируемой выдержки времени \textbf{«Тср Неиспр. ЦУ В1(ЛВ)ВН»} формируется сигнал
\textbf{«Неиспр. ЦУ В1(ЛВ)ВН»}. Выдержка времени должна быть не менее времени взвода пружины привода
выключателя.

\item
Контроль готовности привода выключателя осуществляется по внешнему сигналу \textit{«Пруж. не
заведены В1(ЛВ)ВН»}. Сигнал о незаведенном состоянии пружины \textit{«Ср. Пруж. не заведены В1(ЛВ)ВН»}
формируется с выдержкой времени \textbf{«Тср Пруж. не заведены В1(ЛВ)ВН»}. 

\item
Предусмотрен контроль давления элегаза в баке выключателя по внешним сигналам \textit{«Низк. давл.
SF6 В1(ЛВ)ВН»} и \textit{«Авар. давл. SF6 В1(ЛВ)ВН»} от специальных датчиков давления. Соответствующие
сигналы о низком и аварийном давлении элегаза \textit{«Ср. Низк. давл. SF6 В1(ЛВ)ВН»} и \textit{«Ср. Авар. давл. SF6
В1(ЛВ)ВН»} формируются с выдержками времени, равными 1 с.

\item
Контроль положения автомата шинок питания осуществляется по внешнему сигналу \textit{«Автомат ШП
В1(ЛВ)ВН»}. Сигнал об отключенном положении автомата \textit{«Ср. Автомат ШП В1(ЛВ)ВН»} формируется с
выдержкой времени выдержкой времени, равной 1 с. 

\item
Обогрев выключателя контролируется по внешнему сигналу \textit{«Неиспр. обогр. В1(ЛВ)ВН»}.
Соответствующий сигнал о неисправности обогрева \textit{«Ср. Неиспр. обогр. В1(ЛВ)ВН»} формируется с
регулируемой выдержкой времени \textbf{«Тср Неиспр. обогр. В1(ЛВ)ВН»}.

\item
Сигнализация о неисправности выключателя (сигнал \textit{«Неиспр. В1»}) срабатывает при наличии
одного из сигналов: \textit{«Неиспр. ЦУ В1(ЛВ)ВН»}, \textit{«Ср. Пруж. не заведены В1(ЛВ)ВН»}, \textit{«Ср. Низк. давл. SF6
В1(ЛВ)ВН»}, \textit{«Ср. Авар. давл. SF6 В1(ЛВ)ВН»}, \textit{«Ср. Автомат ШП В1(ЛВ)ВН»}, \textit{«Ср. Неиспр. обогр.
В1(ЛВ)ВН»}, а также от сигнала о непереключении фаз \textit{«Ср. ЗНФ В1(ЛВ)ВН»}.

\begin{figure}[H]
\centering
\includegraphics[width=0.8\textwidth,height=0.8\textheight,keepaspectratio]{\fbpath/761.3371 КСВ/img/CTRL_CT_Q1.pdf}
\captionof{figure}{Функциональная схема контроля исправности ТТ}\label{pic:CTRL_CT_Q1}
\end{figure}

\item
В устройстве также предусмотрен контроль состояния ТТ. Контроль давления элегаза в ТТ
осуществляется по внешним сигналам \textit{«Низк. давл. SF6 ТТ В1(ЛВ)ВН»} и \textit{«Авар. давл. SF6 ТТ В1(ЛВ)ВН»}.
Соответствующие сигналы о низком и аварийном давлении элегаза \textit{«Ср. Низк. давл. SF6 В1(ЛВ)ВН»} и
\textit{«Ср. Авар. давл. SF6 В1(ЛВ)ВН»} формируются с выдержками времени, равными 1 с. Обогрев ТТ
контролируется по внешнему сигналу 
\textit{«Неиспр. обогр. ТТ В1(ЛВ)ВН»}. Соответствующий сигнал о неисправности обогрева \textit{«Ср. Неиспр. обогр.
ТТ В1(ЛВ)ВН»} формируется с регулируемой выдержкой времени \textbf{«Тср Неиспр. обогр. ТТ В1(ЛВ)ВН»}.
По наличию одного из указанных сигналов формируется сигнал о неисправности ТТ \textit{«Неиспр. ТТ
В1(ЛВ)ВН»}. При вводе программного ключа \textbf{«Откл. от Авар. давл. SF6 ТТ В1(ЛВ)ВН»}
предусматривается действие на отключение выключателя при аварийном снижении давления элегаза в
ТТ.

\begin{figure}[H]
\centering
\includegraphics[width=0.8\textwidth,height=0.8\textheight,keepaspectratio]{\fbpath/761.3371 КСВ/img/CTRL_OC_Q1.pdf}
\captionof{figure}{Функциональная схема контроля оперативного тока цепей сигнализации}\label{pic:CTRL_OC_Q1}
\end{figure}

\item
Контроль отсутствия питания цепей сигнализации выключателя и ТТ осуществляется по внешним
сигналам \textit{«Контр. опертока ЭМВ/ЭМО1 В1(ЛВ)ВН»}, \textit{«Контр. опертока ЭМО2 В1(ЛВ)ВН»} и \textit{«Контр.
опертока сигн. В1»}. При состоянии «0» одного из сигналов срабатывает сигнализация о неисправности
цепей сигнализации выключателя и ТТ (сигнал \textit{«Неиспр. сигн. В1/ТТ В1(ЛВ)ВН»}) с выдержкой времени
1 с.  

\begin{figure}[H]
\centering
\includegraphics[width=0.8\textwidth,height=0.8\textheight,keepaspectratio]{\fbpath/761.3371 КСВ/img/PROT_EMU_Q1.pdf}
\captionof{figure}{Функциональная схема защиты ЭМУ от длительного протекания тока}\label{pic:PROT_EMU_Q1}
\end{figure}

\item
Для электромагнитов управления предусмотрена защита от длительного протекания тока.
Контроль протекания тока через ЭМВ, ЭМО1 и ЭМО2 осуществляется при помощи токовых реле,
положения контактов которых должны быть заведены на соответствующие входные сигналы ФСУ
\textit{«Работа ЭМВ В1(ЛВ)ВН»}, \textit{«Работа ЭМО1 В1(ЛВ)ВН»}, \textit{«Работа ЭМО2 В1(ЛВ)ВН»}. При длительном
протекании тока по цепям ЭМО1 и ЭМВ формируется сигнал \textit{«Обест. ЭМО1 и ЭМВ В1(ЛВ)ВН»},
действующий на независимый расцепитель автомата питания цепей ЭМО1 и ЭМВ. Аналогично при
длительном протекании тока по цепи ЭМО2 формируется сигнал \textit{«Обест. ЭМО2 В1(ЛВ)ВН»} для
действия на независимый расцепитель автомата питания цепи ЭМО2. Регулируемые выдержки времени 
\textbf{«Тср защиты ЭМВ В1(ЛВ)ВН»}, \textbf{«Тср защиты ЭМО1 В1(ЛВ)ВН»} и \textbf{«Тср защиты ЭМО2 В1(ЛВ)ВН»}
определяют время срабатывания на обесточивание ЭМВ, ЭМО1 и ЭМО2 соответственно. Кроме того,
наличие сигнала \textit{«Ср. Авар. давл. SF6 В1(ЛВ)ВН»} действует на обесточивание всех электромагнитов.

\begin{figure}[H]
\centering
\includegraphics[width=0.8\textwidth,height=0.8\textheight,keepaspectratio]{\fbpath/761.3371 КСВ/img/BLK_OP_Q1.pdf}
\captionof{figure}{Функциональная схема блокировки управления выключателем}\label{pic:BLK_OP_Q1}
\end{figure}

\item
На блокировку отключения выключателя действует сигнал \textit{«Ср. Авар. давл. SF6 В1(ЛВ)ВН»}.

\item
На блокировку включения действуют сигналы \textit{«Ср. Авар. давл. SF6 В1(ЛВ)ВН»}, \textit{«Неиспр. ЦУ
В1(ЛВ)ВН»}, \textit{«Ср. Пруж. не заведены В1(ЛВ)ВН»}, \textit{«Ср. Автомат ШП В1(ЛВ)ВН»}.

\item
Также предусмотрена блокировка отключения от внешнего сигнала \textit{«Внеш. блок. откл.
В1(ЛВ)ВН»}, блокировка включения от внешнего сигнала \textit{«Внеш. блок. вкл. В1(ЛВ)ВН»} и блокировка
управления (как включения, так и отключения) от внешнего сигнала \textit{«Внеш. блок. упр. В1(ЛВ)ВН»}.

\item
В устройстве предусмотрена программная фиксация команд или включенного положения
выключателя в зависимости от положения программного переключателя \textbf{«Реле фиксации В1(ЛВ)ВН»}.
В положении \textit{«РФК»} фиксируется команда включения, в положении \textit{«РФП»} – включенное положение.
Квитирование сигнала фиксации \textit{«РФ В1(ЛВ)ВН»} может осуществляться от сброса сигнализации или
подачей оперативной команды на отключение. При отключении выключателя от защит и наличии
сигнала \textit{«РФ В1(ЛВ)ВН»} формируется сигнал о факте аварийного отключения выключателя \textit{«Сигнал
несоотв. В1(ЛВ)ВН»}, необходимый для работы АПВ.

\begin{figure}[H]
\centering
\includegraphics[width=0.8\textwidth,height=0.8\textheight,keepaspectratio]{\fbpath/761.3371 КСВ/img/Relay_fix_Q1.pdf}
\captionof{figure}{Функциональная схема фиксации команд или включенного положения выключателя}\label{pic:Relay_fix_Q1}
\end{figure}


\small
\begin{longtable}{|>{\centering\arraybackslash}m{5.3cm}|>{\centering\arraybackslash}m{4.2cm}|>{\centering\arraybackslash}m{1cm}|>{\centering\arraybackslash}m{3.3cm}|>{\centering\arraybackslash}m{1.8cm}|}
% Первая страница: полный заголовок
\caption{Перечень уставочных параметров КСВ\hfill\vspace{-0.5\baselineskip}}\label{desc:tbl1}\\
\hline
\rowcolor{gray!30}
Наименование уставки & Диапазон & Шаг & Значение по умолчанию & Единица измерения \\
\hline
\endfirsthead

% Промежуточные страницы: «Продолжение таблицы»
\caption*{\hspace{3pt}\emph{Продолжение таблицы \ref{desc:tbl1}}\hfill\ }\vspace{-0.5\baselineskip} \\
\hline
\rowcolor{gray!30}
Наименование уставки & Диапазон & Шаг & Значение по умолчанию & Единица измерения \\
\hline
\endhead

%===m>ДТЗ|-
%===m
\end{longtable}
\normalsize	

\end{enumerate}